\section*{Erd\H{o}s Problem 991: equidistribution of logarithmic energy minimizers}
\subsection*{Statement and attribution}
Suppose $p_1,\ldots,p_n\in S^2$ maximize $\prod_{i<j}|p_i-p_j|$.
We prove that their spherical cap discrepancy is $o(n)$.
This qualitative fact is classical potential theory; it is discussed, with quantitative improvements, by Brauchart and by Marzo--Mas. The proof below supplies the compactness, uniqueness, and uniform-cap steps rather than assuming the desired equidistribution. It claims no novelty and does not claim their quantitative decay rates.

\subsection*{A positive expansion of the logarithmic kernel}
Let $\sigma$ be normalized surface area and put
\[
 K(x,y)=\log\frac2{|x-y|},\qquad
 T(x,y)=\frac{1+x\cdot y}{2},\qquad
 K_L(x,y)=\frac12\sum_{k=1}^{L}\frac{T(x,y)^k}{k}.
\]
The nonnegative kernels $K_L$ increase pointwise to $K$, with $K(x,x)=+\infty$.
For fixed $x$, $T(x,y)$ is uniform on $[0,1]$, so
\[
 \iint K\,d\sigma\,d\sigma
 =\frac12\sum_{k\ge1}\frac1{k(k+1)}=\frac12.
 \tag{1}
\]
For any probability measure $\mu$ on the sphere, set $\eta=\mu-\sigma$.
Expansion of $(1+x\cdot y)^k$ into monomials shows
\[
 Q_k(\eta)=\iint T(x,y)^k\,d\eta(x)d\eta(y)\ge0.
\]
Every term is a positive coefficient times the square of a real polynomial moment. The potential of each $T^k$ under $\sigma$ is constant, namely $1/(k+1)$. Therefore, by monotone convergence applied separately to the positive measures,
\[
 \iint K\,d\mu\,d\mu
 =\frac12+\frac12\sum_{k\ge1}\frac{Q_k(\eta)}{k}\ge\frac12.
 \tag{2}
\]
If equality holds, every $Q_k$ is zero. The monomial-square expansion then shows that $\eta$ annihilates every monomial, and hence every polynomial. Polynomials restricted to the sphere are uniformly dense in its continuous functions by the Stone--Weierstrass theorem. Thus $\eta=0$. We have proved that $\sigma$ is the unique energy-minimizing probability measure, including measures whose energy is infinite.

\subsection*{Discrete optimality and weak limits}
Maximizing the distance product is equivalent to minimizing
$\mathcal E_n=\sum_{i\ne j}K(p_i,p_j)$.
A maximum exists by compactness, and it has distinct points because some configuration has positive product. If $n$ independent uniform points are chosen, (1) gives expected energy $n(n-1)/2$. Consequently the minimizing configuration satisfies
\[
 \mathcal E_n\le n(n-1)/2.
\]
Let $\mu_n=n^{-1}\sum_i\delta_{p_i}$. Although its full energy is infinite due to the diagonal, each fixed truncated energy obeys
\[
 \iint K_L\,d\mu_n\,d\mu_n
 \le\frac{\mathcal E_n}{n^2}+\frac{K_L(x,x)}n
 \le\frac12+\frac{H_L}{2n}.
 \tag{3}
\]
Here $K_L(x,x)=H_L/2$ is independent of $x$.
Any subsequence has a further weakly convergent subsequence, because the sphere is compact. If its limit is $\mu$, continuity of $K_L$ passes (3) to
$\iint K_L\,d\mu\,d\mu\le1/2$. Letting $L\to\infty$ and using (2) forces $\mu=\sigma$.
Every subsequential limit is therefore $\sigma$, so $\mu_n$ converges weakly to $\sigma$.

\subsection*{Why the supremum over caps also converges}
Weak convergence alone is pointwise on fixed continuity sets, so a final uniformity argument is necessary.
Write $C(z,t)=\{x:x\cdot z\ge t\}$, $z\in S^2$, $t\in[-1,1]$.
Fix $\delta>0$ and take finite $\delta$-nets for $z$ and $t$.
For any $(z,t)$ choose a net pair $(z',t')$ with
$|z-z'|\le\delta$ and $|t-t'|\le\delta$. Then
\[
 C(z',t'+2\delta)\subseteq C(z,t)\subseteq C(z',t'-2\delta).
\]
Thresholds outside $[-1,1]$ mean the empty or full sphere. These are finitely many comparison caps, and each boundary has $\sigma$-measure zero, including the endpoint singletons.
Their empirical measures converge simultaneously to their areas.
The areas of the two comparison caps differ by at most $2\delta$, since
$\sigma(C(z,u))=(1-u)/2$ on $[-1,1]$.
Sandwiching and then letting $\delta\downarrow0$ gives
\[
 \sup_C|\mu_n(C)-\sigma(C)|\longrightarrow0.
\]
Multiplication by $n$ proves the requested $o(n)$ discrepancy.

\subsection*{Conclusion and sources}
The proof is complete using standard compactness of probability measures and Stone--Weierstrass; the required energy uniqueness was proved explicitly. Sources: \url{https://www.erdosproblems.com/991}; J. S. Brauchart, \emph{Optimal logarithmic energy points on the unit sphere}, Math. Comp. 77 (2008), 1599--1613, \url{https://doi.org/10.1090/S0025-5718-08-02085-1}; J. S. Brauchart, D. P. Hardin and E. B. Saff, \emph{The next-order term for optimal Riesz and logarithmic energy asymptotics on the sphere}, \url{https://doi.org/10.1090/conm/578/11483}.
\subsection*{COMPLETION ESTIMATE}
\textbf{COMPLETION: 100\%}
